% p.78 — الباب الرابع والثلاثون في قوانين التعليل وغيره (تابع: القوانين 11–21)
% Layout: twocol. Right column: rules 11–17 and the first line of 18 (ends «مرفوع اذا»);
% left column: rest of 18 («لم يكن فيه ناصب …»), 19–21. Rule 21 continues on p.79.
\mspage{78}
\begin{twocol}
\begin{rightcol}
(11) اذا كان الالف فى اخر الاسم فلا يتغير فى الحركات الثلاثة نحو مُوسٰى وعِيسٰى
\par
(12) إذا كان حرف العلة متحرّكًا و ما قبله ساكنا فينقل حركته الى ما قبله لان حرف العلة ضعيف لا يحمل الحركة وما قبله حرف صحيح قوى ساكن فيحمل الحركة نحو يَقْوُلُ يَقُوْلُ يَبْيِعُ يَبِيْعُ
\par
(13) ثم يقلب حرف العلة مطابقا لما قبله كما مر فتقول فى يَخْوَفُ يَخَوْفُ اولا ثم يَخَافُ ثَانِيًا
\par
(14) لا يجوز التقاء الساكنين فلا يقال \unclear{قَرْءْ مِنْ الابجد لا} وَهُوَ ما يأتى
% TODO: the example phrase after «فلا يقال» is hard to read (قَرْءْ مِنْ الابجد لا ?); the muḥaqqiq should check the lithograph.
\par
(15) اذا كان الحرف مشددا وما قبله حرف مد فيجوز التقاء الساكنين مثل دَابَّةٍ دِيْبَّةٍ دُوْبَّةٍ
\par
(16) حرف المد الف ما قبله مفتوح و واو ما قبله مضموم وياء ما قبله مكسور مثل قَالَ قِيْلَ قُوْلَ
\par
(17) لا يجوز التقاء الساكنين ولا الابتداء بالسّكون
\par
(18) اعلم ان المضارع مرفوع اذا
\end{rightcol}
\begin{leftcol}
لم يكن فيه ناصب و جازم مثل يَضْرِبُ ويَنْصُرُ وفى الجزم لم يضربْ ولم ينصُرْ وفى النصب لن يضربَ ولن ينصرَ
\par
(19) اذا كان المضارع فعلا ناقصا فلا يظهر الرفع والجزم بل يقال فى الرفع يَغْزُوْ ويَرْمِىْ ويَخْشٰى بابقاء الواو والياء والالف وفى الجزم لَمْ يَغْزُ ولَمْ يَرْمِ و لَمْ يخش بالحذف لان الحذف علامة الجزم وفى النصب لن يغزوَ و لن يرميَ ولن يَخْشٰى باظهار النصب
\par
(20) اذا كان المضارع فعلا اجوف تقول فى الرفع يَقُوْلُ ويَبِيْعُ ويَخافُ وفى النصب لن يقولَ ولن يبيعَ و لن يخافَ باظهار الرفع والنصب وفى الجزم لم يَقُلْ ولم يَبِعْ و لم يَخَفْ بحذف الواو والياء والالف لان التقاء الساكنين لا يجوز
\par
(21) واذا كان المضارع مُثَنًّى او \unclear{مَجْمُوْعًا} فيقال فى الرفع يضربان يضربون ينصران ينصرون باثبات النون لانها علامة الرفع وفى الجزم والنصب لم يضربا لم يضربوا لن يضربا لن يضربوا بحذفها
\end{leftcol}
\end{twocol}
